\documentclass[11pt,twocolumn,a4paper]{article}

% ─── Packages ─────────────────────────────────────────────────────────
\usepackage[margin=2.0cm,top=2.4cm,bottom=2.4cm]{geometry}
\usepackage[utf8]{inputenc}
\usepackage[T1]{fontenc}
\usepackage{lmodern}
\usepackage{microtype}
\usepackage{amsmath,amssymb,amsthm,mathtools}
\usepackage{bm}
\usepackage{graphicx}
\usepackage{booktabs}
\usepackage{array}
\usepackage{enumitem}
\usepackage{float}
\usepackage{caption}
\usepackage{natbib}
\usepackage{xcolor}
\usepackage{listings}
\usepackage[colorlinks=true,linkcolor=blue!60!black,
            citecolor=blue!60!black,urlcolor=blue!60!black]{hyperref}
\usepackage{fancyhdr}
\usepackage{titlesec}
\usepackage{cleveref}

% ─── Page Style ───────────────────────────────────────────────────────
\pagestyle{fancy}
\fancyhf{}
\fancyhead[L]{\small\textit{Heihachi: A Runtime Graph for Music Production}}
\fancyhead[R]{\small\thepage}
\renewcommand{\headrulewidth}{0.4pt}

\titleformat{\section}{\large\bfseries}{\thesection.}{0.5em}{}
\titleformat{\subsection}{\normalsize\bfseries}{\thesubsection.}{0.5em}{}

% ─── Theorem Environments ─────────────────────────────────────────────
\newtheorem{theorem}{Theorem}[section]
\newtheorem{lemma}[theorem]{Lemma}
\newtheorem{corollary}[theorem]{Corollary}
\newtheorem{proposition}[theorem]{Proposition}
\theoremstyle{definition}
\newtheorem{definition}[theorem]{Definition}
\newtheorem{example}[theorem]{Example}
\newtheorem{principle}[theorem]{Principle}
\theoremstyle{remark}
\newtheorem{remark}[theorem]{Remark}

% ─── Custom Commands ──────────────────────────────────────────────────
\newcommand{\R}{\mathbb{R}}
\newcommand{\N}{\mathbb{N}}
\newcommand{\Rpos}{\R_{>0}}
\newcommand{\flo}{\beta}
\newcommand{\med}{\mathcal{M}}
\newcommand{\Gph}{\Gamma}
\newcommand{\cut}{\operatorname{cut}}
\newcommand{\wt}{w}
\newcommand{\res}{\varrho}
\newcommand{\rec}{\mu}
\newcommand{\Nodes}{\mathcal{N}}
\newcommand{\Chunks}{\operatorname{Ch}}
\newcommand{\Vals}{\operatorname{Val}}
\newcommand{\Traj}{\mathcal{T}}
\newcommand{\node}{n}
\newcommand{\dispatch}{\textsc{Run}}
\newcommand{\Modules}{\mathcal{M}\!od}

% ─── Listings ─────────────────────────────────────────────────────────
\lstdefinestyle{hk}{
  basicstyle=\ttfamily\footnotesize,
  keywordstyle=\bfseries,
  commentstyle=\itshape\color{gray!80!black},
  showstringspaces=false,
  breaklines=true,
  columns=fullflexible,
  frame=single,
  framesep=4pt,
  rulecolor=\color{gray!50},
  morekeywords={node,chunk,value,module,emit,read,run,commit,identify,
                transform,for,in,do,end,if,then,else,return,let}
}
\lstset{style=hk}

\captionsetup{font=small,labelfont=bf}

% ══════════════════════════════════════════════════════════════════════
\begin{document}

\twocolumn[
\begin{@twocolumnfalse}
\begin{center}

{\LARGE\bfseries Heihachi: A Runtime Graph for Music Production\\[4pt]
A Kernel that Executes, Records, and Judges Nothing}

\vspace{12pt}

{\large Kundai Farai Sachikonye}

\vspace{6pt}

{\small\itshape Independent Researcher \\[0.3em]
\texttt{kundai.sachikonye@wzw.tum.de}}

\vspace{6pt}{\small August 2026}

\vspace{16pt}

\begin{minipage}{0.92\textwidth}
\textbf{Abstract.}
We specify \emph{heihachi}, a runtime for music production organised as an
accumulating record of authored decisions rather than as a pipeline of
processing steps. The runtime manages a graph whose nodes are \emph{subtasks} of
production---``the transient of this bass,'' ``the width of this pad''---each
carrying a bag of executable realisations and a growing collection of values. Its
entire semantic content is one operation: execute every realisation attached to a
node and increment a counter once per emission. It does not schedule, does not
select among realisations, does not inspect what they produce, and cannot report
success or failure. We prove that this last incapacity is structural rather than
elective: a verdict presupposes a comparison against an expectation, the graph
provides no vocabulary in which an expectation can be written, and consequently
no quantity the runtime computes can carry the meaning of an exit code. Three
consequences follow and we establish each. Runs reach completion in the presence
of anomalies, because nothing branches on emission content---measured over a
sweep of forty-eight chain-length and anomaly-density combinations, the runtime
completes every node in every cell while a halt-on-error runtime falls to
$9.4\%$ of the work and abandons up to twenty-nine nodes. The trajectory of a run
is constituted by the reads and emissions that occurred and is therefore not
storable as a schedule; over one fixed six-node catalogue, thirty-six runs induce
six distinct causal edge relations, so no single fixed relation describes them
all. And relevance is constituted by being read, so an emission nothing consumes
leaves the trajectory without any actor having rejected it. We give the value
type and prove that no value of non-positive resolution is constructible; we give
the record and prove it admits no cached recomputation, so that re-measuring is a
new commitment at a strictly higher record; and we give a hierarchical address
space in which an edit at depth $k$ of a $b$-ary tree of depth $D$ affects exactly
$b^{D-k}$ leaves, measured across branching factors two, three and four. A
reference implementation and forty deterministic experiments accompany the paper.
Nothing outside this paper is required to read it.
\end{minipage}

\vspace{8pt}
\noindent\textbf{Keywords:}
runtime graph, micro-kernel, music production, provenance, exit code,
run to completion, causal edges, monotone record, resolution floor

\vspace{18pt}
\end{center}
\end{@twocolumnfalse}
]

%──────────────────────────────────────────────────────────────────────
\section{Introduction}
\label{sec:intro}
%──────────────────────────────────────────────────────────────────────

\subsection{What is lost between sessions}

A producer working on a record makes some thousands of decisions. A bass is
synthesised, resampled, mangled, resampled again; a version is bounced and set
aside; another is bounced and kept. Six months later the kept version is the only
survivor, and the decisions that produced it are gone---not archived somewhere
inconvenient, but genuinely absent, because nothing recorded them.

This is ordinarily described as a documentation problem, to be fixed by better
naming conventions or more disciplined session hygiene. We think that
description is wrong, and the argument against it is short. A session file stores
\emph{state}: which devices are present, what their parameters are. A bounce
stores \emph{output}: what the state produced. Neither stores the relation between
them, and the relation is what a later question is about. ``Why does this one
sound right and the other twenty not'' is a question about the twenty, and the
twenty were discarded.

The asymmetry is sharper still. A bounce can be regenerated from a session; a
session cannot be regenerated from a bounce. The processing chain is not
invertible, and no amount of analysis of the output recovers the sequence of
authored choices that produced it. Whatever a production tool is for, it cannot
be for reconstructing that sequence after the fact, because the reconstruction is
not available.

Heihachi therefore records the relation while it is being made. This paper
specifies the runtime that does the recording. Two companion languages---one for
computing over the record, one for constructing new material---are specified
separately; this paper defines the substrate both run on, and is self-contained.

\subsection{Three refusals}

The design is easiest to state by what the runtime declines to do, because in
each case something structural does the job that the declined machinery would
have done.

\begin{itemize}[leftmargin=1.4em]
\item \textbf{No schedule.} Nodes are individuated by the subtask they name, not
  by the agent that raised them. Two analyses that decompose different problems
  and arrive at the same subtask converge on one node
  (\Cref{def:convergence}). There is consequently no ordering decision for a
  scheduler to make.
\item \textbf{No orchestration.} An emission that nothing reads forms no causal
  edge and does not enter the trajectory (\Cref{thm:nonprop}). Because relevance
  is constituted by being read, no actor is required to reject the irrelevant.
\item \textbf{No exit code.} The graph has no vocabulary for ``expected'' and
  stores no expectation. The runtime therefore holds nothing against which an
  achieved value could be measured, and cannot issue a verdict
  (\Cref{thm:no-exit}).
\end{itemize}

What replaces the missing machinery is a single premise---that telling one thing
from another costs a strictly positive amount---from which the value type, the
record, and the addressing scheme all follow.

\subsection{Why the third refusal matters here}

The no-exit-code result is the one that sounds like a defect and is in fact the
reason the design suits its domain.

Consider what a conventional runtime does with an anomaly. It compares the
anomalous value against an expectation, finds a mismatch, and halts. In a
compiler or a build system this is correct: a type error is not a result, it is
the absence of one. In a recording studio it is not correct, because an
unexpected reading is frequently the point. A resample that came out wrong, a
compressor that pumped in an unplanned way, a filter that self-oscillated---these
are the events a producer builds a record out of. A runtime that halted on them
would be halting on the material.

The wet-lab discipline is the right analogy. An unexpected reading does not abort
the assay; it is written in the notebook and becomes the thing the analysis has to
explain. \Cref{cor:completion} establishes the same discipline formally: an
anomaly is a first-class value on a node, available to any module that reads it
and to the final report, and the report---not an exit code---is the runtime's
output.

\subsection{Contributions}

\begin{enumerate}[label=(\roman*),leftmargin=2.2em]
\item The positive floor, the value type it induces, and a proof that no value of
  non-positive resolution is constructible (\S\ref{sec:floor}).
\item The node, the chunk bag, and convergence by subtask identity
  (\S\ref{sec:node}).
\item The kernel's single operation, the no-exit-code theorem, and run to
  completion, with the abandoned-work measurement (\S\ref{sec:runtime}).
\item The monotone record, its use as an internal arrow of time, and the
  no-memoisation corollary (\S\ref{sec:record}).
\item Run-induced edges, trajectory emergence, and non-propagation without
  rejection (\S\ref{sec:trajectory}).
\item Accountability as a minimum cut, checked against an independent
  computation (\S\ref{sec:accountability}).
\item Resolution as a depth in one address space, with the blast-radius law
  (\S\ref{sec:resolution}).
\item Floor negotiation between a program and a render path
  (\S\ref{sec:negotiation}).
\item A reference implementation and forty deterministic experiments
  (\S\ref{sec:validation}).
\end{enumerate}

%──────────────────────────────────────────────────────────────────────
\section{The Floor}
\label{sec:floor}
%──────────────────────────────────────────────────────────────────────

\subsection{Contact graphs}

\begin{definition}[Contact graph, medium, cut, residue]
\label{def:substrate}
A \emph{contact graph} is $\Gph=(V,E,\wt)$ with $V$ finite and non-empty, $E$ a
set of unordered pairs, and $\wt\colon E\to\Rpos$; graph terminology follows
\citet{diestel2010graph}. A distinguished vertex $\med\in V$, the \emph{medium},
is adjacent to every other vertex; the remaining vertices are \emph{items}. For
$S\subseteq V$ containing an item and omitting $\med$, the \emph{cut} is
$\cut(S)=\{\{u,v\}\in E: u\in S,\, v\notin S\}$ and the \emph{residue} is
$\res(S)=\sum_{e\in\cut(S)}\wt(e)$. The \emph{floor} is
$\flo=\min_{e\in E}\wt(e)$.
\end{definition}

The medium is the vertex standing for everything not yet told apart. In a
production context it is the unstructured remainder: the material that has not
been distinguished from anything, against which a first distinction is drawn.

\begin{lemma}[Positive floor, positive residue]
\label{lem:floor}
$\flo>0$, and every cut as in \Cref{def:substrate} satisfies
$\res(S)\ge\flo>0$.
\end{lemma}

\begin{proof}
$\flo$ is the minimum of a finite non-empty set of strictly positive reals and is
therefore positive. Since $S$ contains an item $v$ and omits $\med$, and
$\{v,\med\}\in E$ by the medium property, $\cut(S)$ is non-empty; its weight is a
sum of terms each at least $\flo$.
\end{proof}

\begin{remark}[Where the floor comes from in audio]
\label{rem:audio-floor}
\Cref{lem:floor} is a graph fact, but the quantity it names is an ordinary
engineering one. A spectral measurement resolves no better than its bin width; a
level meter with $N$ output states over a range $D$ distinguishes nothing below
$D/N$; a converter of $b$ bits over a $5$\,V span does not resolve differences
below $5/2^{b}$ volts. These are statements about instrument specifications, not
about noise, and they set the weights of \Cref{def:substrate} in practice.
\end{remark}

\subsection{Values}

\begin{definition}[Value]
\label{def:value}
A \emph{value} is a datum attached to a node, carrying a magnitude, a strictly
positive \emph{floor} recording the resolution at which it was obtained, a unit,
a kind, and the record index at which it was committed. A value of non-positive
floor is not constructible.
\end{definition}

\begin{theorem}[No zero-resolution value]
\label{thm:no-zero}
There is no constructible value whose floor is zero or negative. Equivalently, a
claim of exact measurement is not expressible.
\end{theorem}

\begin{proof}
Immediate from \Cref{def:value}, whose introduction rule carries the side
condition $\flo>0$; there is no other value-introducing rule. The content of the
restriction is \Cref{lem:floor}: a value is the record of a distinction, a
distinction is a cut, and every cut has residue at least $\flo>0$. A value of
zero floor would be the record of a distinction drawn at no cost, and no such cut
exists.
\end{proof}

The experiment \texttt{kernel\_floor\_positivity} attempts three
non-positive floors and records three refusals, with a positive floor accepted.

\begin{remark}[What the floor is not]
\label{rem:not-clamp}
The floor is the unit in which residues are expressed, not a clamp applied after
the fact. If residues are multiples of the floor, rescaling the floor rescales
them and leaves the ratio $\res/\flo$ unchanged. The experiment
\texttt{kernel\_floor\_is\_scale} sweeps the floor over a sixteen-fold range and
finds the ratio constant to $0$ in double precision---which is the operational
test that an implementation treats the floor as a scale rather than a truncation.
\end{remark}

%──────────────────────────────────────────────────────────────────────
\section{The Node}
\label{sec:node}
%──────────────────────────────────────────────────────────────────────

\subsection{Modules, subtasks, chunks}

The system is a set of \emph{modules}: a spectral analyser, a groove estimator, a
loudness meter, an annotation store, a model that reads the graph. A module is a
competence rather than a procedure, and it solves nothing directly. Given a
problem it decomposes it into \emph{subtasks}, and this decomposition is the only
source of structure in the system. There is no global plan.

A subtask is not private to the module that raised it. When a decomposition
arrives at a subtask, that subtask is published: it becomes a node readable by
every module.

\begin{definition}[Node]
\label{def:node}
A \emph{node} is a triple
\[
\node=\bigl(\tau,\ \Chunks(\node),\ \Vals(\node)\bigr),
\]
where $\tau$ is the identity of a subtask; $\Chunks(\node)=\{c_1,\dots,c_r\}$ is
a finite \emph{bag} of chunks, each an executable realisation of $\tau$; and
$\Vals(\node)$ is the time-varying collection of values the node carries.
\end{definition}

\begin{definition}[Chunk]
\label{def:chunk}
A \emph{chunk} is an executable fragment realising a subtask. The runtime
requires only that a chunk can be executed and that its execution produces a
finite collection of value emissions. The runtime does not parse, type-check,
interpret or compare chunks.
\end{definition}

\begin{remark}[Why a bag]
\label{rem:bag}
The multiplicity of $\Chunks(\node)$ is essential. A subtask of production can be
realised in more than one way: ``the transient of this bass'' may be realised by
a spectral onset detector, by a stored annotation from a previous session, and by
a model's reading of the graph. These are not competing candidates among which
the runtime must choose. They are co-resident realisations, and by
\Cref{def:exec} the runtime executes all of them. Selecting among them would
require a criterion; a criterion is a judgement; and judgements do not belong to
the runtime.
\end{remark}

\subsection{Convergence}

\begin{definition}[Convergence]
\label{def:convergence}
Two subtask realisations \emph{converge} when they share subtask identity $\tau$.
Convergence merges their chunk bags and unifies their value collections onto the
single node bearing $\tau$. The graph therefore grows by convergence: raising a
subtask that already exists adds chunks and values to an existing node rather
than creating a new one.
\end{definition}

\begin{remark}[A node is a meeting point]
\label{rem:meeting}
A node is not owned by the module that first raised it. Because the subtask
individuates the node, and because many modules converge on the same subtask, a
node accretes chunks and values from every module that reaches it. Its content is
the union of what converged there. This is the mechanism by which a heterogeneous
set of analyses accumulates shared structure with no central registry of who
computes what.
\end{remark}

The experiment \texttt{kernel\_convergence} raises one subtask from two modules
and records a single node bearing two chunks and two values.

\subsection{The graph vocabulary}

\begin{definition}[Graph vocabulary]
\label{def:vocab}
A module's interface to the graph consists of exactly four operations:
\emph{identify} (locate the node bearing a subtask identity, raising it if it
does not exist), \emph{read} (retrieve the values on a node), \emph{transform}
(alter its own internal state), and \emph{emit} (attach new values to a node).
There is no fifth operation.
\end{definition}

\begin{principle}[The value contract]
\label{prin:contract}
Every module emits values onto nodes and reads values from nodes. A module reads
the values on a node, makes internal changes on the strength of what it read, and
emits new values---possibly onto the same node, possibly onto others. There is no
other channel of inter-module communication. No module reads another module's
internal state, and no module is required to interpret another module's chunks.
\end{principle}

\begin{remark}[The vocabulary is a closure claim]
\label{rem:closure}
\Cref{def:vocab} is not an inventory but a closure claim, and it is what forces
adjudication out of the runtime. There is no operation \textsc{IsCorrect}, no
operation \textsc{Expected}, and no operation \textsc{Compare}. The graph stores
what was emitted, never what should have been emitted. A module cannot ask the
graph whether a bass sounds right, because the graph represents no such property;
it can only read what is there and emit on the strength of it. Taste therefore
remains in the modules, where it belongs, and the runtime never enforces
anybody's current opinion.
\end{remark}

%──────────────────────────────────────────────────────────────────────
\section{The Runtime}
\label{sec:runtime}
%──────────────────────────────────────────────────────────────────────

\subsection{Execution is the only responsibility}

\begin{definition}[Node execution]
\label{def:exec}
To \emph{execute} a node $\node$ is to run every chunk in its bag:
\[
\dispatch(\node)=\bigl\langle \textsc{Eval}(c) : c\in\Chunks(\node)\bigr\rangle,
\]
each chunk producing whatever value emissions it produces, all of which are
attached to the graph. The runtime performs $\dispatch(\node)$ and nothing else.
It does not select among chunks; it does not order the graph; it does not inspect
the values that result.
\end{definition}

\begin{lstlisting}[caption={The kernel's entire semantic content.},label={lst:kernel}]
run(node):
  for chunk in node.chunks do
    execute(chunk)      -- emissions land on the graph
    commit()            -- record += 1
  end
\end{lstlisting}

\begin{principle}[Semantic inertia]
\label{prin:inert}
The runtime is a device for turning nodes into value emissions by running their
code. It holds no representation of an expected value, no comparison operator
against one, and no policy over emission content.
\end{principle}

\subsection{No exit code}

\begin{theorem}[No exit code]
\label{thm:no-exit}
The runtime cannot emit an exit code, because an exit code presupposes a
comparison the runtime is structurally unable to make.
\end{theorem}

\begin{proof}
An exit code is a verdict: success against a class of failure. To issue one, the
issuer must compare an achieved state to an expected one. The runtime's entire
interface with meaning is the value collection $\Vals$, and by \Cref{rem:closure}
the graph has no verb for ``expected'' and stores no expectation against which an
achieved value could be measured. The runtime therefore has no expected state to
compare against and no comparison operator in its vocabulary. Whatever a chunk
emits is emitted; there is no state in which the runtime holds a belief about
what should have been emitted. Hence there is no quantity the runtime could
compute carrying the meaning of an exit code. Adjudication is possible only
downstream, in a consumer holding an expectation of its own, and a consumer's
judgement is that consumer's, not the runtime's.
\end{proof}

The experiment \texttt{kernel\_no\_exit\_code} inspects the runtime's report and
finds its keys to be \texttt{record}, \texttt{nodes\_executed},
\texttt{emissions}, \texttt{anomalies} and \texttt{induced\_edges}: no key
carrying a verdict, with the anomaly recorded as a value alongside the ordinary
result.

\begin{corollary}[Run to completion]
\label{cor:completion}
The runtime does not halt on an unexpected result or on an error emitted by a
chunk. It runs each node's chunks and moves on; anomalous emissions are recorded
as values like any other.
\end{corollary}

\begin{proof}
By \Cref{thm:no-exit} the runtime has no notion of ``unexpected,'' so there is no
event of that type on which to halt. A chunk that raises produces a value---an
error record---emitted onto the graph exactly as a numeric result would be.
Nothing in $\dispatch(\node)$ branches on the \emph{content} of an emission, so
the presence of an error value cannot alter control flow.
\end{proof}

\begin{table}[H]
\centering
\small
\caption{Work abandoned by halting. Measured over eight chain lengths
$L\in\{4,\dots,32\}$ and six anomaly densities $d\in\{0.0,\dots,0.5\}$
(\texttt{kernel\_run\_to\_completion}, 48 cells).}
\label{tab:completion}
\begin{tabular}{@{}lr@{}}
\toprule
Quantity & Value \\
\midrule
Cells swept & 48 \\
Cells in which heihachi completed every node & 48 \\
Worst-case fraction completed by halting runtime & $0.094$ \\
Greatest number of nodes abandoned by halting & 29 \\
\bottomrule
\end{tabular}
\end{table}

\begin{remark}[Where errors are handled]
\label{rem:errors-elsewhere}
\Cref{cor:completion} does not say errors are ignored. It says they are not
handled \emph{by the runtime}. A module that reads an error value may do anything
it likes with it: retry through a different chunk, emit a diagnostic, degrade a
downstream estimate, or raise a finding for the report. That is a module's
semantic prerogative. The runtime's refusal to act on error content is what keeps
the semantics in the modules.
\end{remark}

\begin{proposition}[Total chunk execution]
\label{prop:total}
For every node execution, every chunk in the node's bag is executed exactly once,
and no chunk is skipped on the basis of the values produced by any other chunk.
\end{proposition}

\begin{proof}
Immediate from \Cref{def:exec}, whose iteration ranges over the whole bag and
whose body does not read the emissions of previous iterations.
\end{proof}

The experiment \texttt{kernel\_total\_chunk\_execution} places a raising chunk
second in a bag of four and records all four as having fired.

%──────────────────────────────────────────────────────────────────────
\section{The Record}
\label{sec:record}
%──────────────────────────────────────────────────────────────────────

\begin{definition}[Record]
\label{def:record}
Let $\rec$ assign to each emission the value of a global strictly monotone
counter, incremented once per emission. Emission $u$ \emph{precedes} $v$, written
$u\prec v$, if there is a directed provenance path from $u$ to $v$.
\end{definition}

\begin{theorem}[Internal arrow of time]
\label{thm:arrow}
$\prec$ is a strict partial order and $\rec$ is a linear extension of it.
Consequently, for any two provenance-connected emissions, deciding which caused
which requires only the internal counter and no external clock.
\end{theorem}

\begin{proof}
Provenance paths are directed and acyclic, since an emission's provenance names
values that existed before it was committed and $\rec$ strictly increases at each
commitment; hence $\prec$ is irreflexive and transitive. If $u\prec v$ then $u$
was committed before $v$ and so $\rec(u)<\rec(v)$, which is the linear-extension
property. The argument follows the logical-clock construction of
\citet{lamport1978time}.
\end{proof}

\begin{theorem}[Monotonicity and non-return]
\label{thm:monotone}
The record is strictly increasing along any execution. A system state that
repeats a prior \emph{configuration} of values does not return to a prior
\emph{state}: the configuration recurs but the record has strictly grown.
\end{theorem}

\begin{proof}
Each emission increments $\rec$ by one and no operation decrements it. To
decrease $\rec$ a committed emission would have to be withdrawn; a withdrawal is
itself an emission---a value recording the retraction---which increments $\rec$
rather than decrementing it.
\end{proof}

\begin{corollary}[No cached recomputation]
\label{cor:no-cache}
Re-evaluating a measurement is a new emission at a strictly higher record, never
a cached recomputation. A stored reading returned later is a determination made
against the graph \emph{as it was} at the earlier record; since no operation
restores that earlier graph, returning it would present a determination of a
state the system has irreversibly left.
\end{corollary}

The experiment \texttt{kernel\_monotone\_record} performs four hundred
measurements, records zero decreases, then repeats an identical measurement three
times and obtains three distinct records---$401$, $402$, $403$---confirming that
no memoisation occurred.

\begin{remark}[What this buys a producer]
\label{rem:record-producer}
\Cref{cor:no-cache} sounds like a restriction and is in fact the property that
makes a session history usable. If re-analysing a bounce could return a cached
answer, the history would contain determinations whose relation to the material
is unstated: was this measured before or after the resample? Because every
determination sits at a record index, and the index orders every commitment in the
session, that question always has an answer.
\end{remark}

%──────────────────────────────────────────────────────────────────────
\section{The Trajectory}
\label{sec:trajectory}
%──────────────────────────────────────────────────────────────────────

\begin{proposition}[Edges are induced, not declared]
\label{prop:edges}
An edge $u\to v$ obtains in a run if and only if, in that run, a value emitted at
$u$ was read by a module that then emitted at $v$. The edge set is therefore a
product of a run, not an input to it, and two runs over the same node set may
induce different edge sets.
\end{proposition}

Nodes are the durable objects; edges are records of causation that occurred, not
declarations of dependency to be scheduled. This is the sense in which the graph
is causal.

\begin{definition}[Trajectory]
\label{def:trajectory}
The \emph{trajectory} $\Traj$ of a run is the set of nodes that carried
propagated information---the nodes at which some module read a value and, on the
strength of it, emitted a value that another module later read---together with
the values as they stood when read.
\end{definition}

\begin{theorem}[Trajectory emergence]
\label{thm:emergence}
The trajectory of a run is not knowable before the run and is not storable as a
schedule. It can only be produced by running.
\end{theorem}

\begin{proof}
Membership of a node in $\Traj$ depends on whether some module read a value there
and emitted onward. Whether a module reads a given value is a decision the module
makes on the strength of values it has already read in this run
(\Cref{prin:contract}). Thus the trajectory at step $t$ is a function of the
values present at step $t$, which are the emissions of steps $<t$, which depend
on the trajectory up to $t-1$. The relation is a fixed point of a map taking a
partial trajectory to its own extension, and because the map's action at each
step consults run-produced values---not any datum available before the run---no
evaluation is possible without executing the steps it depends on.
\end{proof}

The experiment \texttt{kernel\_trajectory\_emergence} runs one fixed six-node
catalogue thirty-six times under varying seeds and magnitudes and recovers six
distinct induced edge relations, so no single fixed relation describes them all.

\begin{corollary}[What is durable]
\label{cor:durable}
The node set---``what can be computed''---is durable and may be catalogued,
imported and frozen. The trajectory---``what was computed''---cannot.
Reproducibility must therefore attach to the former.
\end{corollary}

\begin{theorem}[Non-propagation without rejection]
\label{thm:nonprop}
Let $\node$ be a node whose emissions no module reads. Then no causal edge
incident to $\node$ is formed and $\node\notin\Traj$. Moreover no actor performed
a rejection.
\end{theorem}

\begin{proof}
By \Cref{prop:edges} an edge incident to $\node$ requires a module to have read a
value at $\node$ and emitted onward; by hypothesis none did, so no such edge
exists, and by \Cref{def:trajectory} $\node\notin\Traj$. The construction of
$\Traj$ consults only reads that occurred; nothing in it evaluates a value and
declines it.
\end{proof}

\begin{remark}[Relevance is constituted, not detected]
\label{rem:constituted}
There is no true-but-faint signal being discarded here, because relevance is
constituted by being read rather than being a prior property that a filter might
fail to detect. This is why the design needs no orchestrator: the question
``which of these emissions matter?'' is answered by what the modules did, not by
an authority that ruled on it.
\end{remark}

The experiment \texttt{kernel\_non\_propagation} emits at two nodes, reads onward
from one, and finds the other absent from the trajectory with zero rejections
performed.

%──────────────────────────────────────────────────────────────────────
\section{Accountability}
\label{sec:accountability}
%──────────────────────────────────────────────────────────────────────

The runtime issues no verdict, but a consumer may still ask why something
happened. That question has a precise answer.

\begin{definition}[Separation cost]
\label{def:sep}
For items $u,v$ of a contact graph, the \emph{separation cost} $\sigma(u,v)$ is
the minimum weight of a cut separating them, computable by maximum flow
\citep{ford1956maximal,edmonds1972theoretical}. The \emph{strength} of an item is
its separation from the medium, $\sigma(v)=\sigma(v,\med)$.
\end{definition}

\begin{theorem}[Accountability is a minimum cut]
\label{thm:mincut}
Let $s$ be a virtual source joined to every exogenous emission of a run and let
$v^{\star}$ be the node bearing an outcome of interest. The set of emissions
jointly necessary for $v^{\star}$---those whose removal would remove the
outcome---is a minimum-weight $s$--$v^{\star}$ cut in the provenance graph with
capacities $\wt$.
\end{theorem}

\begin{proof}
A set of emissions whose removal disconnects $v^{\star}$ from $s$ is exactly an
$s$--$v^{\star}$ cut, since provenance paths from exogenous inputs to
$v^{\star}$ are the routes by which influence reached it. Among such sets the
jointly necessary one is of minimum weight, since any lighter cut would
disconnect $v^{\star}$ at lower cost and the heavier set would contain a
member whose removal leaves $v^{\star}$ connected---hence not necessary. Minimum
cut equals maximum flow \citep{ford1956maximal}.
\end{proof}

\begin{remark}[Why a cut and not a correlation]
\label{rem:mincut-why}
The cut formulation answers ``which decisions were \emph{jointly} necessary for
this outcome?'' Correlational attribution answers a different question---which
decisions co-varied with it---and offers no guarantee that removing the
attributed set would remove the outcome. For a producer asking why one bounce
works and another does not, joint necessity is the question being asked.
\end{remark}

The experiment \texttt{kernel\_mincut\_brute\_force} checks the maximum-flow
implementation against a brute-force minimum over all source--sink bipartitions
across three hundred random graphs, finding zero mismatches with a greatest
absolute difference of $8.9\times 10^{-16}$. The theorem is therefore checked
against an independent computation rather than against itself.

\begin{remark}[Cost and placement]
\label{rem:offline}
The minimum cut is a post-hoc query, polynomial but not constant-time. It belongs
in a report, not in a render loop.
\end{remark}

%──────────────────────────────────────────────────────────────────────
\section{Resolution as Depth}
\label{sec:resolution}
%──────────────────────────────────────────────────────────────────────

A production tool must let a user work at many scales: the whole arrangement, one
section, one bar, one transient. Conventional tools implement these as separate
modes with separate interfaces. Heihachi treats them as depths in one address
space.

\begin{definition}[Hierarchical address]
\label{def:address}
An \emph{address} is a finite sequence $a=\langle a_1,\dots,a_k\rangle$ read
coarse to fine. The \emph{resolution} of a determination is the depth $k$ at
which it is expressed. Addresses form a tree of branching factor $b$ and depth
$D$; a prefix names an interior node and every extension of it a descendant.
\end{definition}

\begin{proposition}[Blast radius]
\label{prop:blast}
In a uniform address tree of branching factor $b$ and depth $D$, an edit applied
at depth $k$ affects exactly $b^{D-k}$ leaf objects.
\end{proposition}

\begin{proof}
The subtree rooted at a depth-$k$ address has height $D-k$, and a uniform tree of
height $h$ and branching factor $b$ has exactly $b^{h}$ leaves. An edit at the
root of that subtree applies to its leaves and to no others.
\end{proof}

\begin{table}[H]
\centering
\small
\caption{Dynamic range of control, $D=4$
(\texttt{kernel\_blast\_radius}).}
\label{tab:blast}
\begin{tabular}{@{}lrr@{}}
\toprule
Branching $b$ & Surgical edit & Coarsest edit \\
\midrule
$2$ & $1$ leaf & $16$ leaves \\
$3$ & $1$ leaf & $81$ leaves \\
$4$ & $1$ leaf & $256$ leaves \\
\bottomrule
\end{tabular}
\end{table}

\begin{remark}[One space, not two artefacts]
\label{rem:one-space}
The payoff is that a stock setting and a hand-tuned one are the same object at
two depths rather than two artefacts. In a conventional tool a preset and the
patch it expands to are different things touched at different times; here they
are one address space traversed to whatever depth a decision requires, and the
blast radius of an edit is known before it is made.
\end{remark}

%──────────────────────────────────────────────────────────────────────
\section{Floor Negotiation}
\label{sec:negotiation}
%──────────────────────────────────────────────────────────────────────

A program declares a floor; a render path has a numeric resolution. These are
independent quantities---the first chosen by the author, the second fixed by the
machinery---and if the second is coarser than the first, the program's
distinctions are not representable.

\begin{definition}[Numeric resolution]
\label{def:numres}
The \emph{numeric resolution} $\varepsilon_T$ of a back end is the largest
absolute error it may introduce in computing a residue: accumulated rounding for
a floating-point evaluator \citep{goldberg1991what,higham2002accuracy}, plus the
quantisation step where results are read back from a discretised buffer.
\end{definition}

\begin{theorem}[Floor negotiation]
\label{thm:negotiate}
A program with ambient floor $\flo$ may execute on a back end of numeric
resolution $\varepsilon_T$, against a reference of resolution $\varepsilon_R$,
only if $\varepsilon_R+\varepsilon_T\le\flo$. Otherwise the residues it computes
are not distinguishable from the errors it incurs, and the program must be
refused before execution.
\end{theorem}

\begin{proof}
Two values whose residues differ by less than $\varepsilon_R+\varepsilon_T$
cannot be told apart by comparing computed residues, since the computation may
perturb each by up to its resolution. The program's floor asserts that
differences of size $\flo$ are meaningful. If $\varepsilon_R+\varepsilon_T>\flo$
then a difference the program treats as meaningful lies within the machinery's
error, so a reported distinction may be an artefact.
\end{proof}

\begin{corollary}[Refusal, not silent disagreement]
\label{cor:refuse}
A program whose floor is finer than its render path can resolve is refused, with
both quantities named in the refusal.
\end{corollary}

\begin{remark}[The two quantities are independent]
\label{rem:indep}
It is tempting to argue that the back end's error is automatically bounded by the
floor, on the grounds that the floor is the coarsest resolvable distinction. The
argument is invalid. The floor is a property of the weights the author chose;
$\varepsilon_T$ is a property of the arithmetic and buffer format the render path
fixes. Nothing couples them, which is exactly why the check must be performed
rather than assumed.
\end{remark}

%──────────────────────────────────────────────────────────────────────
\section{Interfaces}
\label{sec:interfaces}
%──────────────────────────────────────────────────────────────────────

\begin{lstlisting}[caption={The complete module interface.},label={lst:module}]
module M:
  identify(tau)      -> node
  read(node)         -> values
  transform(values)  -> internal   -- private
  emit(node, values) -> ()
\end{lstlisting}

\begin{lstlisting}[caption={The complete chunk contract.},label={lst:chunk}]
chunk c:
  execute(context) -> emissions
  -- must terminate within its declared budget
  -- may raise: a raise produces an error value
\end{lstlisting}

A chunk that exceeds its budget is pre-empted and an emission recording the
pre-emption is appended. This is a timing act, not a judgement: no expectation
about the chunk's output is consulted, only the fact that its budget elapsed.

\begin{principle}[Separation of content and timing]
\label{prin:separation}
The runtime may not compare a value to an expectation. It may nonetheless
guarantee that operations complete within a bounded interval. These are
independent: a deadline is a fact about elapsed duration and contains no
reference to emission content.
\end{principle}

%──────────────────────────────────────────────────────────────────────
\section{Validation}
\label{sec:validation}
%──────────────────────────────────────────────────────────────────────

A reference implementation accompanies this paper: \texttt{hkcore.py} (contact
graphs, maximum flow and minimum cut, the runtime, values and the record),
\texttt{exp\_kernel.py}, and \texttt{run\_all.py}. It uses the Python standard
library only. Each experiment fixes its seed and writes a JSON record; the suite
reports \textbf{40 run, 40 passed, 0 failed} in $0.16$\,s, of which eleven
experiments belong to this paper. Re-running reproduces every experiment record
bit-identically.

\begin{table}[H]
\centering
\small
\caption{Experiments for this paper.}
\label{tab:experiments}
\begin{tabular}{@{}lp{4.3cm}@{}}
\toprule
Experiment & Establishes \\
\midrule
\texttt{floor\_positivity} & \Cref{thm:no-zero}: three refusals of three attempts \\
\texttt{mincut\_brute\_force} & \Cref{thm:mincut} against brute force; 0/300 mismatches \\
\texttt{run\_to\_completion} & \Cref{cor:completion}; 48 cells, 29 nodes recovered \\
\texttt{no\_exit\_code} & \Cref{thm:no-exit}: no verdict key in the report \\
\texttt{monotone\_record} & \Cref{thm:monotone}, \Cref{cor:no-cache}: 0 decreases \\
\texttt{convergence} & \Cref{def:convergence}: two modules, one node \\
\texttt{total\_chunk\_execution} & \Cref{prop:total}: 4 of 4 chunks fired \\
\texttt{trajectory\_emergence} & \Cref{thm:emergence}: 6 relations over 36 runs \\
\texttt{non\_propagation} & \Cref{thm:nonprop}: 0 rejections performed \\
\texttt{blast\_radius} & \Cref{prop:blast}: spans $1$--$256$ at $b=4$ \\
\texttt{floor\_is\_scale} & \Cref{rem:not-clamp}: ratio spread $0$ \\
\bottomrule
\end{tabular}
\end{table}

\begin{remark}[What the suite does not establish]
\label{rem:suite-limits}
The suite establishes that the definitions are coherent and implementable and
that an implementation behaves as the theorems predict. It does not establish
that the model is adequate for any particular studio, that the floors chosen for
audio measurements are the right ones, or that the runtime performs acceptably at
session scale. \Cref{sec:limits} states where the design should not be used.
\end{remark}

%──────────────────────────────────────────────────────────────────────
\section{Related Notions}
\label{sec:related}
%──────────────────────────────────────────────────────────────────────

The runtime resembles several established models and differs from each in a way
worth stating.

\emph{Dataflow and stream processing.} Kahn process networks
\citep{kahn1974semantics} and synchronous dataflow \citep{lee1987synchronous}
schedule computation over a declared dependency graph. Heihachi declares no such
graph: by \Cref{prop:edges} edges are induced by a run, so the topology is an
output rather than an input.

\emph{Actors and message passing.} The actor model \citep{hewitt1973universal}
gives each actor private state and an address. Heihachi's modules likewise hold
private state, but they do not address one another at all: the only channel is
the graph (\Cref{prin:contract}), and a node is individuated by a subtask rather
than by an identity a sender must know.

\emph{Tuple spaces.} Linda \citep{gelernter1985generative} shares the property
that producers and consumers are decoupled through a medium. The difference is
that a tuple is consumed by being read, whereas a value on a node persists and
accretes, and that heihachi's medium is structured by subtask identity rather
than by pattern matching.

\emph{Provenance systems.} Workflow provenance \citep{davidson2008provenance} and
logical clocks \citep{lamport1978time} supply the record's ancestry. The
departure is that provenance here is not an audit trail attached to a
computation but the computation's own output, since by \Cref{thm:no-exit} there
is no other result to report.

\emph{Audio processing graphs.} Music-DSP environments
\citep{puckette1996pure,mccartney2002rethinking} represent signal flow as a graph
of unit generators evaluated in a topologically sorted order. That graph and this
one are different objects: theirs carries samples between processors at audio
rate, this one carries authored decisions between analyses at whatever rate they
occur. A production system needs both, and they should not be conflated.

%──────────────────────────────────────────────────────────────────────
\section{Limits and Non-Applicability}
\label{sec:limits}
%──────────────────────────────────────────────────────────────────────

\emph{This is not an audio processing graph.} The runtime carries values, not
samples. It schedules nothing and guarantees no sample-accurate timing. A system
that renders audio needs a separate real-time path, and \Cref{prin:separation}
is what keeps the two from contaminating each other.

\emph{Storage grows without bound.} The record is monotone and the log is
append-only, so a long session accumulates. \Cref{rem:audio-floor}'s floors and
the contact floor bound what is retained, but the growth is real and an
implementation must plan for it.

\emph{Do not use this where bit-identical replay is the product.} A system whose
value is exact reproduction of a computation---a mastering render that must match
last week's byte for byte, a regression oracle, a delivery pipeline under
contract---should not adopt this architecture at its centre. By
\Cref{cor:durable} what repeats is the node catalogue and the setup, never the
trajectory. The runtime is for the part of production that measures a world it
does not control; the final render is not that part.

\emph{The runtime cannot make anything sound good.} It has no verdict
(\Cref{thm:no-exit}) and holds no expectation. It records what was decided and
what followed. Every judgement remains with the producer and with the modules
they choose to trust.

%──────────────────────────────────────────────────────────────────────
\section{Conclusion}
\label{sec:conclusion}
%──────────────────────────────────────────────────────────────────────

We have specified a runtime whose whole semantic content is
\Cref{lst:kernel}: execute every chunk in a node's bag, and increment a counter
per emission. From that and one premise---that telling one thing from another
costs a strictly positive floor---follow the value type
(\Cref{thm:no-zero}), the impossibility of a verdict (\Cref{thm:no-exit}), run to
completion (\Cref{cor:completion}), the emergence of the trajectory
(\Cref{thm:emergence}), non-propagation without rejection (\Cref{thm:nonprop}),
accountability as a minimum cut (\Cref{thm:mincut}), and the blast-radius law
(\Cref{prop:blast}).

What the design offers a producer is narrow and, we think, the right narrow
thing. It does not generate material, evaluate it, or recommend anything. It
records the decisions as they are made, in an order that survives, with the
resolution of every claim carried in the claim itself---so that the question
asked six months later, about the twenty versions that were discarded and the one
that was not, has somewhere to be asked.

\bibliographystyle{plainnat}
\bibliography{references}

\end{document}
